%% GENERATED by tools/gen_latex.py -- edit content/ch17.py instead.
%% Build:  latexmk -pdf ch17.tex     (or: pdflatex ch17.tex, twice)
\documentclass[aspectratio=169,11pt,t]{beamer}
\usepackage{itbpro}

\coursetitle{Image Generation}
\coursesubtitle{Two ways to build a latent space of images --- one by construction, one by denoising in a loop --- and a walk between two prompts that shows what a neural network really is.}
\coursekicker{Chapter 17}
\coursesource{Chollet \& Watson, Deep Learning with Python 3e --- chapter 17}
\courseduration{3.5 hours (3 sessions)}
\coursepresenter{Rahman Indra Kesuma, S.Kom., M.Cs. $\cdot$ Prof. Bambang Riyanto Trilaksono}{Teaching Assistant $\cdot$ Lead Instructor}
\coursebrand{AI for Professional \textperiodcentered\ ITB \textperiodcentered\ Ch 17}

\title{Image Generation}
\author{Rahman Indra Kesuma, S.Kom., M.Cs. $\cdot$ Prof. Bambang Riyanto Trilaksono}
\date{}

\begin{document}
\covertitle

\framekicker{Learning outcomes}
\begin{frame}[shrink=25]{Session Objectives}
\leadpar{By the end of this session, participants can:}
\begin{isteps}
\item Explain what a \textbf{latent space of images} is, and why sampling from one produces images nobody has drawn.
\item Describe a \textbf{classical autoencoder} and say why its latent space is not useful for generation.
\item Build a \textbf{VAE}: encoder to a mean and variance, a sampling layer, a decoder --- and both halves of its loss.
\item Customise training with \textbf{compute\_loss()} rather than train\_step(), and say why that keeps the code backend-agnostic.
\item Explain \textbf{reverse diffusion} as a denoising autoencoder in a loop, and define diffusion time and the diffusion schedule.
\item Build a \textbf{U-Net denoiser} that predicts a noise mask, and train it on the Oxford Flowers dataset.
\item Use a pretrained \textbf{Stable Diffusion} text-to-image model, with negative prompts and a controllable step count.
\item Interpolate between two prompts with \textbf{slerp}, and explain why spherical interpolation beats linear.
\end{isteps}

\vskip 6pt
\vskip 4pt
\reslink{Book}{Chapter 17 --- full text}{https://deeplearningwithpython.io/chapters/chapter17\_image-generation}
\reslink{Notebook}{01\_vae\_on\_mnist.ipynb}{https://hendrikarisma.my.id/rs/ai-products-course/notebooks/ch17/01\_vae\_on\_mnist.html}
\reslink{Notebook}{02\_vae\_latent\_space\_grid.ipynb}{https://hendrikarisma.my.id/rs/ai-products-course/notebooks/ch17/02\_vae\_latent\_space\_grid.html}
\reslink{Notebook}{03\_diffusion\_unet\_and\_schedule.ipynb}{https://hendrikarisma.my.id/rs/ai-products-course/notebooks/ch17/03\_diffusion\_unet\_and\_schedule.html}
\reslink{Notebook}{04\_training\_the\_flower\_diffuser.ipynb}{https://hendrikarisma.my.id/rs/ai-products-course/notebooks/ch17/04\_training\_the\_flower\_diffuser.html}
\reslink{Notebook}{05\_stable\_diffusion\_latent\_walk.ipynb}{https://hendrikarisma.my.id/rs/ai-products-course/notebooks/ch17/05\_stable\_diffusion\_latent\_walk.html}
\reslink{Notebook}{All chapter 17 notebooks}{https://hendrikarisma.my.id/rs/ai-products-course/notebooks/index.html\#ch17}
\reslink{Repo}{Author's official notebook}{https://github.com/fchollet/deep-learning-with-python-notebooks/blob/master/chapter17\_image-generation.ipynb}
\vskip 2pt
\end{frame}

\coursesection{01}{Latent spaces of images}{The one idea underneath every image generator.}

\framekicker{Section 17.1.1}
\begin{frame}[shrink=25]{A space where every point is a valid image}
\leadpar{The key idea of image generation: develop a \textbf{low-dimensional latent space} --- a vector space, like everything else in deep learning --- in which any point maps to an image that looks like the real thing.}
{\fontsize{9.4}{12.6}\selectfont\color{fg2}The module that realises the mapping, taking a latent point and outputting a grid of pixels, is called a \textbf{generator}, or sometimes a \textbf{decoder}.\par}\vskip 4pt
\begin{center}
\includegraphics[width=0.92\linewidth,height=0.60\textheight,keepaspectratio]{figs/ch17-latent.tex.pdf}
\end{center}
\vskip 3pt{\color{ink3}\fontsize{7.4}{9.4}\selectfont Figure 17.1 --- once the space is learned, sample from it and decode.\par}
{\fontsize{9.4}{12.6}\selectfont\color{fg2}What comes out are the \hilite{in-betweens of the training images}.\par}\vskip 4pt
\end{frame}

\framekicker{Section 17.1.1}
\begin{frame}[shrink=25]{Add text conditioning and you get a brush}
{\fontsize{9.4}{12.6}\selectfont\color{fg2}\textbf{Text conditioning} maps a space of natural-language prompts into the latent space, giving language-guided image generation. This category is called \textbf{text-to-image models}.\par}\vskip 4pt
\begin{center}
\includegraphics[width=0.92\linewidth,height=0.60\textheight,keepaspectratio]{figs/ch17-t2i.tex.pdf}
\end{center}
\vskip 3pt{\color{ink3}\fontsize{7.4}{9.4}\selectfont Figure 17.2 --- a prompt lands somewhere in the joint space, and the decoder renders it.\par}
{\fontsize{9.4}{12.6}\selectfont\color{fg2}Interpolating between many training images lets such models generate \textbf{infinite combinations of visual concepts}, including many nobody had explicitly come up with. A horse riding a bike on the moon? You got it.\par}\vskip 4pt
\end{frame}

\framekicker{Section 17.1.1 $\cdot$ the honest caveat}
\begin{frame}[shrink=25]{The latent space does not model the physical world}
{\fontsize{9.4}{12.6}\selectfont\color{fg2}Like every deep learning model, the latent space encodes no consistent model of physics. You will occasionally see hands with extra fingers, incoherent lighting, garbled objects. Coherence is an area of active research.\par}\vskip 4pt
\begin{iband}[amber]
Despite having seen tens of thousands of images of people riding bikes, the model does not understand in a human sense what riding a bike \textit{means} --- pedalling, steering, balance. Which is why your horse is \hilite{unlikely to be depicted pedalling with its hind legs} the way a human artist would draw it.
\end{iband}
{\fontsize{9.4}{12.6}\selectfont\color{fg2}This is chapter 15's interpolation argument, arriving in a different modality. \textbf{Nothing about pixels changes it.}\par}\vskip 4pt
\end{frame}

\framekicker{Section 17.1.1 $\cdot$ the landscape}
\begin{frame}[shrink=25]{Three families, and why one is missing}
\begin{minipage}[t]{0.3253\textwidth}
\begin{icardaccent}
\cardh{Diffusion models}
The architecture behind \textbf{nearly all commercial image generation services today}. Section 17.2.
\end{icardaccent}
\end{minipage}%
\hfill
\begin{minipage}[t]{0.3253\textwidth}
\begin{icardaccent}
\cardh{Variational autoencoders}
Highly structured, controllable latent spaces. Not the first choice for fidelity, but still in use. Section 17.1.
\end{icardaccent}
\end{minipage}%
\hfill
\begin{minipage}[t]{0.3253\textwidth}
\begin{icard}
\cardh{GANs}
Covered in previous editions of this book. They have \textbf{gradually fallen out of fashion} and been all but replaced by diffusion.
\end{icard}
\end{minipage}%
{\fontsize{9.4}{12.6}\selectfont\color{fg2}Note also that none of this is image-specific. You could develop latent spaces of \textbf{sound or music} with the same models --- pictures are simply where the most interesting results have been obtained so far.\par}\vskip 4pt
\end{frame}

\coursesection{02}{Variational autoencoders}{A little statistical magic that forces a latent space to be continuous.}

\framekicker{Section 17.1.2}
\begin{frame}[shrink=25]{Start with the classical autoencoder}
{\fontsize{9.4}{12.6}\selectfont\color{fg2}An autoencoder maps an image to a latent vector space via an \textbf{encoder}, and decodes it back to the same dimensions via a \textbf{decoder}. It is trained with the input images as targets --- it learns to reconstruct its own input.\par}\vskip 4pt
\begin{center}
\includegraphics[width=0.92\linewidth,height=0.60\textheight,keepaspectratio]{figs/ch17-autoencoder.tex.pdf}
\end{center}
\vskip 3pt{\color{ink3}\fontsize{7.4}{9.4}\selectfont Figure 17.3 --- the target is the input. This is self-supervised learning.\par}
{\fontsize{9.4}{12.6}\selectfont\color{fg2}By constraining the \textbf{code} --- most commonly to be low dimensional and sparse --- the encoder becomes a way of compressing the input into fewer bits.\par}\vskip 4pt
\end{frame}

\framekicker{Section 17.1.2}
\begin{frame}[shrink=25]{And note that classical autoencoders do not work well}
\leadpar{In practice they lead to latent spaces that are neither useful nor nicely structured. \textbf{They are not much good at compression either.} For these reasons they have largely fallen out of fashion.}
{\fontsize{9.4}{12.6}\selectfont\color{fg2}VAEs augment the autoencoder with a little statistical magic that forces it to learn \textbf{continuous, highly structured} latent spaces --- and that turns out to be a powerful tool for image generation.\par}\vskip 4pt
{\fontsize{9.4}{12.6}\selectfont\color{fg2}VAEs were discovered simultaneously by \textbf{Kingma and Welling} in December 2013 and by \textbf{Rezende, Mohamed, and Wierstra} in January 2014. They mix deep learning with Bayesian inference, and remain in use in research a decade later.\par}\vskip 4pt
\end{frame}

\framekicker{Section 17.1.2 $\cdot$ the mechanism}
\begin{frame}[shrink=25]{Encode to a distribution, not to a point}
{\fontsize{9.4}{12.6}\selectfont\color{fg2}Instead of compressing an input into a \textbf{fixed code}, a VAE turns the image into the parameters of a statistical distribution: a \textbf{mean} and a \textbf{variance}.\par}\vskip 4pt
{\fontsize{9.4}{12.6}\selectfont\color{fg2}The assumption is that the image was generated by a statistical process, and that the randomness of that process should be accounted for during encoding and decoding.\par}\vskip 4pt
\begin{center}

\includegraphics[width=0.92\linewidth,height=0.60\textheight,keepaspectratio]{figs/ch17-vae.tex.pdf}
\end{center}
\vskip 3pt{\color{ink3}\fontsize{7.4}{9.4}\selectfont Figure 17.4 --- the encoder outputs two vectors; a point is sampled from the distribution they define; the decoder reconstructs from that point.\par}
\end{frame}

\framekicker{Section 17.1.2 $\cdot$ three steps}
\begin{frame}[shrink=25]{The mechanism, stated precisely}
\begin{isteps}
\item An \textbf{encoder} turns \inlinecode{input\_img} into two parameters in the latent space, \inlinecode{z\_mean} and \inlinecode{z\_log\_variance}.
\item A point is \textbf{randomly sampled} from the latent normal distribution assumed to have generated the image: \inlinecode{z = z\_mean + exp(z\_log\_variance) * epsilon}, where \inlinecode{epsilon} is a random tensor of small values.
\item A \textbf{decoder} maps that point back to the original input image.
\end{isteps}
\begin{iband}[signal]
Because \inlinecode{epsilon} is random, \textbf{every point close to where \inlinecode{input\_img} was encoded decodes to something similar to \inlinecode{input\_img}.} That is what forces the latent space to be \hilite{continuously meaningful}.
\end{iband}
\end{frame}

\framekicker{Section 17.1.2 $\cdot$ why continuity matters}
\begin{frame}[shrink=25]{Continuity plus low dimensionality gives you concept vectors}
{\fontsize{9.4}{12.6}\selectfont\color{fg2}Any two close points in the latent space decode to highly similar images. Combine that continuity with a \textbf{low dimensional} space and every direction is forced to encode a meaningful axis of variation in the data.\par}\vskip 4pt
{\fontsize{9.4}{12.6}\selectfont\color{fg2}The result is a very structured space, highly suitable to manipulation via \textbf{concept vectors} --- the same \textit{word arithmetic} idea from chapter 15, now over faces and objects rather than words.\par}\vskip 4pt
{\fontsize{9.4}{12.6}\selectfont\color{fg2}This is why VAEs remain relevant even where diffusion produces better images: when \textbf{interpretability, control over the latent space, and reconstruction} matter, they are still the right tool.\par}\vskip 4pt
\end{frame}

\framekicker{Section 17.1.2 $\cdot$ the loss}
\begin{frame}[shrink=25]{Two losses, pulling in different directions}
\begin{center}
\includegraphics[width=0.92\linewidth,height=0.60\textheight,keepaspectratio]{figs/ch17-vae-loss.tex.pdf}
\end{center}
\vskip 3pt{\color{ink3}\fontsize{7.4}{9.4}\selectfont Reconstruction alone would collapse to a classical autoencoder; regularization alone would ignore the data.\par}
{\fontsize{9.4}{12.6}\selectfont\color{fg2}The regularization term is the \textbf{Kullback-Leibler divergence}, which nudges the distribution of the encoder output toward a well-rounded normal distribution centred on 0 --- a sensible prior about the structure of the space being modelled.\par}\vskip 4pt
\end{frame}

\framekicker{Section 17.1.3 $\cdot$ listing 17.1}
\begin{frame}[shrink=25]{The encoder: a ConvNet with two heads}
{\fontsize{9.4}{12.6}\selectfont\color{fg2}A simple ConvNet mapping a 28 $\times$ 28 MNIST digit to two two-dimensional vectors. Note the \textbf{strides} for downsampling, not max pooling.\par}\vskip 4pt
\icode{python}{listing 17.1}{listings/ch17-01.py}
{\fontsize{9.4}{12.6}\selectfont\color{fg2}Strides are preferable to max pooling for \textbf{any model that cares where things are in the image} --- as chapter 11's segmentation example did, and as this one does, since the encoding must support reconstructing a valid image.\par}\vskip 4pt
\end{frame}

\framekicker{Section 17.1.3}
\begin{frame}[shrink=25]{Sixty-nine thousand parameters, and a two-dimensional output}
\ioutput{listings/ch17-02.txt}
{\fontsize{9.4}{12.6}\selectfont\color{fg2}The whole of MNIST is being pushed through a \textbf{two-number} bottleneck. That is deliberate: a 2D latent space is one we can plot in full, which is the point of this example.\par}\vskip 4pt
\end{frame}

\framekicker{Section 17.1.3 $\cdot$ listing 17.2}
\begin{frame}[shrink=25]{The sampling layer}
{\fontsize{9.4}{12.6}\selectfont\color{fg2}One custom layer applies the VAE sampling formula. It needs a seed generator to use \inlinecode{keras.random} inside \inlinecode{call()}.\par}\vskip 4pt
\icode{python}{listing 17.2}{listings/ch17-03.py}
{\fontsize{9.4}{12.6}\selectfont\color{fg2}\inlinecode{exp(0.5 * z\_log\_var)} converts a \textbf{log variance} to a standard deviation. Predicting the log rather than the variance itself keeps the encoder's output unconstrained --- it may be negative, and the exponential makes it positive.\par}\vskip 4pt
\end{frame}

\framekicker{Section 17.1.3 $\cdot$ listing 17.3}
\begin{frame}[shrink=25]{The decoder mirrors the encoder}
{\fontsize{9.4}{12.6}\selectfont\color{fg2}Reshape the latent vector to image dimensions, then reverse each encoder operation with \inlinecode{Conv2DTranspose}.\par}\vskip 4pt
\icode{python}{listing 17.3}{listings/ch17-04.py}
{\fontsize{9.4}{12.6}\selectfont\color{fg2}The \inlinecode{Dense(7 * 7 * 64)} produces exactly the number of coefficients the encoder's \inlinecode{Flatten} consumed --- the architecture is a \textbf{mirror}, and reading it that way is the quickest way to check it.\par}\vskip 4pt
\end{frame}

\framekicker{Section 17.1.3}
\begin{frame}[shrink=25]{compute\_loss() instead of train\_step()}
{\fontsize{9.4}{12.6}\selectfont\color{fg2}This is the first model in the book that is \textbf{not doing supervised learning} --- an autoencoder is self-supervised, using its inputs as targets. Departing from supervised learning normally means subclassing \inlinecode{Model} and writing a custom \inlinecode{train\_step()}.\par}\vskip 4pt
\begin{iband}[signal]
But \inlinecode{train\_step()} contents are \textbf{backend specific} --- \inlinecode{GradientTape} under TensorFlow, \inlinecode{loss.backward()} under PyTorch. Overriding \inlinecode{compute\_loss()} instead keeps the default \inlinecode{train\_step()} and stays \hilite{backend agnostic}.
\end{iband}
\icode{python}{the signature}{listings/ch17-05.py}
{\fontsize{9.4}{12.6}\selectfont\color{fg2}\inlinecode{x} is the input, \inlinecode{y} the target --- \textbf{None} here, since our dataset has no targets --- and \inlinecode{y\_pred} is the output of \inlinecode{call()}. The method returns a scalar, and may also update metric state.\par}\vskip 4pt
\end{frame}
\note{Worth pausing on. This is a small API decision with a large practical consequence: it is the difference between code that runs on three backends and code that runs on one.}

\framekicker{Section 17.1.3 $\cdot$ listing 17.4}
\begin{frame}[shrink=25]{The VAE model: constructor}
{\fontsize{9.4}{12.6}\selectfont\color{fg2}Encoder, decoder, sampler, and two metrics to track the loss components separately over each epoch.\par}\vskip 4pt
\icode{python}{listing 17.4 - constructor}{listings/ch17-06.py}
{\fontsize{9.4}{12.6}\selectfont\color{fg2}Note that \inlinecode{call()} returns only the \textbf{encoder output} --- the two latent parameters. Sampling and decoding happen inside the loss, because that is where the reconstruction is needed.\par}\vskip 4pt
\end{frame}

\framekicker{Section 17.1.3 $\cdot$ listing 17.4}
\begin{frame}[shrink=25]{The VAE model: the loss}
{\fontsize{9.4}{12.6}\selectfont\color{fg2}Both terms, written out. \inlinecode{x} is the original image; \inlinecode{y\_pred} carries the latent parameters.\par}\vskip 4pt
\icode{python}{listing 17.4 - compute\_loss}{listings/ch17-07.py}
{\fontsize{9.4}{12.6}\selectfont\color{fg2}The reconstruction loss is \textbf{summed over the spatial dimensions} (axes 1 and 2) and averaged over the batch. The KL term is the closed-form divergence between the encoder's distribution and a unit normal.\par}\vskip 4pt
\end{frame}

\framekicker{Section 17.1.3 $\cdot$ listing 17.5}
\begin{frame}[shrink=25]{Training with no loss and no targets}
{\fontsize{9.4}{12.6}\selectfont\color{fg2}Because \inlinecode{compute\_loss()} already handles the loss, we pass no \inlinecode{loss} at compile time --- which in turn means passing \textbf{no target data} to \inlinecode{fit()}.\par}\vskip 4pt
\icode{python}{listing 17.5}{listings/ch17-08.py}
{\fontsize{9.4}{12.6}\selectfont\color{fg2}We concatenate the training and test splits: there is no test set to protect here, because there is \textbf{nothing being evaluated against held-out labels}. We want the latent space to see every digit there is.\par}\vskip 4pt
\end{frame}

\framekicker{Section 17.1.3 $\cdot$ listing 17.6}
\begin{frame}[shrink=25]{Walking the latent space on a grid}
{\fontsize{9.4}{12.6}\selectfont\color{fg2}With a 2D latent space we can plot the whole thing: sample points linearly on a grid and decode each one.\par}\vskip 4pt
\icode{python}{listing 17.6}{listings/ch17-09.py}
{\fontsize{9.4}{12.6}\selectfont\color{fg2}Nine hundred digits, none of which are in MNIST.\par}\vskip 4pt
\end{frame}

\framekicker{Section 17.1.3 $\cdot$ figure 17.5}
\begin{frame}[shrink=25]{What the grid shows}
\leadpar{A \textbf{completely continuous} distribution of the different digit classes, with one digit morphing into another as you follow a path through the latent space.}
{\fontsize{9.4}{12.6}\selectfont\color{fg2}Specific directions in the space have meaning: there is a direction for \textit{four-ness}, one for \textit{one-ness}, and so on. Nobody labelled them. They fall out of the continuity constraint.\par}\vskip 4pt
\begin{iband}[amber]
This is the payoff for the whole VAE construction. A classical autoencoder's latent space would show \textbf{islands of digits separated by regions that decode to nothing} --- the KL term is what fills the gaps.
\end{iband}
\end{frame}

\coursesection{03}{Diffusion models}{A denoising autoencoder in a loop --- and the architecture behind nearly every commercial image generator today.}

\framekicker{Section 17.2}
\begin{frame}[shrink=25]{Denoising is the one task autoencoders excel at}
{\fontsize{9.4}{12.6}\selectfont\color{fg2}A long-standing application of autoencoders: feed in an input with a small amount of noise --- a low-quality JPEG, say --- and get back a cleaned-up version.\par}\vskip 4pt
{\fontsize{9.4}{12.6}\selectfont\color{fg2}In the late 2010s this gave rise to successful \textbf{image super-resolution} models, which have shipped in every major smartphone camera app for years.\par}\vskip 4pt
\begin{iband}[signal]
These models are \textbf{not} recovering lost detail hidden in the input, like the \textit{enhance} scene in \textit{Blade Runner}. They are making educated guesses --- \hilite{hallucinating} a higher-resolution version of what you gave them.
\end{iband}
\end{frame}

\framekicker{Section 17.2 $\cdot$ a cautionary example}
\begin{frame}[shrink=25]{The moon that was not there}
{\fontsize{9.4}{12.6}\selectfont\color{fg2}With some AI-enhanced cameras you can photograph something \textbf{vaguely moon-like} --- a printout of a severely blurred moon image --- and find a crisp picture of the moon's craters in your camera roll.\par}\vskip 4pt
{\fontsize{9.4}{12.6}\selectfont\color{fg2}A great deal of detail that simply was not present in the printout gets straight-up hallucinated, because the super-resolution model is \textbf{overfitted to moon photography}.\par}\vskip 4pt
\begin{iband}[rose]
So, unlike Rick Deckard, \hilite{definitely do not use this technique for forensics}.
\end{iband}
{\fontsize{9.4}{12.6}\selectfont\color{fg2}Worth keeping for any professional audience: an enhancement model's output is \textbf{evidence of the model's priors}, not evidence about the scene.\par}\vskip 4pt
\end{frame}

\framekicker{Section 17.2 $\cdot$ the key insight}
\begin{frame}[shrink=25]{If a little noise, why not all of it?}
\begin{center}
\includegraphics[width=0.92\linewidth,height=0.60\textheight,keepaspectratio]{figs/ch17-diffusion-idea.tex.pdf}
\end{center}
\vskip 3pt{\color{ink3}\fontsize{7.4}{9.4}\selectfont The arresting idea that early denoising successes led researchers to.\par}
{\fontsize{9.4}{12.6}\selectfont\color{fg2}These should more accurately be called \textbf{reverse diffusion} models --- \textit{diffusion} refers to the forward process of gradually adding noise to an image until it disperses into nothing.\par}\vskip 4pt
\end{frame}

\framekicker{Section 17.2 $\cdot$ figure 17.7}
\begin{frame}[shrink=25]{Reverse diffusion, one step at a time}
\begin{center}
\includegraphics[width=0.92\linewidth,height=0.60\textheight,keepaspectratio]{figs/ch17-reverse.tex.pdf}
\end{center}
\vskip 3pt{\color{ink3}\fontsize{7.4}{9.4}\selectfont Figure 17.7 --- turning pure noise into an image via repeated denoising.\par}
\begin{iquote}
\quotetext{\textit{"Every block of stone has a statue inside it and it is the task of the sculptor to discover it."} --- Well, \textbf{every square of white noise has an image inside it, and it is the task of the diffusion model to discover it.}}\quotecite{Michelangelo, and section 17.2}
\end{iquote}
\end{frame}

\framekicker{Section 17.2.1}
\begin{frame}[shrink=25]{Oxford Flowers: 8,189 images, 102 species}
{\fontsize{9.4}{12.6}\selectfont\color{fg2}A small dataset, which is what makes this example trainable in an afternoon.\par}\vskip 4pt
\icode{python}{}{listings/ch17-10.py}
{\fontsize{9.4}{12.6}\selectfont\color{fg2}Two arguments matter. \inlinecode{labels=None} --- we want images only. \inlinecode{crop\_to\_aspect\_ratio=True} --- resizing without it would \textbf{distort} the images, and distortion in the training set becomes distortion in everything generated.\par}\vskip 4pt
\end{frame}

\framekicker{Section 17.2.2}
\begin{frame}[shrink=25]{The denoiser predicts the noise, not the image}
{\fontsize{9.4}{12.6}\selectfont\color{fg2}The same denoising model is reused at every iteration of the loop, erasing a little noise each time. To make its job easier, we \textbf{tell it how much noise to expect} --- the \inlinecode{noise\_rates} input.\par}\vskip 4pt
\begin{iband}[signal]
And rather than outputting a denoised image, the model outputs a \textbf{predicted noise mask}, which we subtract from the input. \hilite{Predicting what to remove is an easier target than predicting what remains.}
\end{iband}
{\fontsize{9.4}{12.6}\selectfont\color{fg2}For the architecture we use a \textbf{U-Net} --- a ConvNet originally developed for image segmentation, and last seen in chapter 11.\par}\vskip 4pt
\end{frame}

\framekicker{Section 17.2.2 $\cdot$ figure 17.9}
\begin{frame}[shrink=25]{Three stages, with skip connections across}
\begin{center}
\includegraphics[width=0.92\linewidth,height=0.60\textheight,keepaspectratio]{figs/ch17-unet.tex.pdf}
\end{center}
\vskip 3pt{\color{ink3}\fontsize{7.4}{9.4}\selectfont Figure 17.9 --- a 1:1 mapping between downsampling and upsampling blocks.\par}
{\fontsize{9.4}{12.6}\selectfont\color{fg2}Each upsampling block is the inverse of a downsampling block, and \textbf{concatenative residual connections} run from each downsampling block to its partner. These avoid the loss of image detail across successive down- and upsampling.\par}\vskip 4pt
\end{frame}

\framekicker{Section 17.2.2}
\begin{frame}[shrink=25]{The residual block}
{\fontsize{9.4}{12.6}\selectfont\color{fg2}A utility applying two convolutions with a residual connection, projecting the shortcut with a 1 $\times$ 1 convolution when the widths do not match.\par}\vskip 4pt
\icode{python}{}{listings/ch17-11.py}
{\fontsize{9.4}{12.6}\selectfont\color{fg2}This is the pattern from chapter 9, unchanged --- with \inlinecode{swish} in place of \inlinecode{relu}, and normalization that does \textbf{not} learn a scale or centre, since the residual path carries those.\par}\vskip 4pt
\end{frame}

\framekicker{Section 17.2.2}
\begin{frame}[shrink=25]{The U-Net: two inputs, and the way down}
{\fontsize{9.4}{12.6}\selectfont\color{fg2}The two inputs are merged first, then the downsampling stage saves each block's output on a stack.\par}\vskip 4pt
\icode{python}{get\_model - inputs and downsampling}{listings/ch17-12.py}
{\fontsize{9.4}{12.6}\selectfont\color{fg2}The scalar noise rate is \textbf{upsampled to the full image size} and concatenated as an extra channel --- the standard way to feed a scalar condition to a convolutional network.\par}\vskip 4pt
\end{frame}

\framekicker{Section 17.2.2}
\begin{frame}[shrink=25]{The U-Net: the middle, and the way back up}
{\fontsize{9.4}{12.6}\selectfont\color{fg2}The upsampling loop runs the widths in reverse, popping one skip connection per block.\par}\vskip 4pt
\icode{python}{get\_model - middle and upsampling}{listings/ch17-13.py}
{\fontsize{9.4}{12.6}\selectfont\color{fg2}\inlinecode{skips.pop()} pairs each upsampling block with its mirror on the way down --- \textbf{last in, first out} is exactly the pairing the architecture diagram shows.\par}\vskip 4pt
\end{frame}

\framekicker{Section 17.2.2 $\cdot$ two details worth stealing}
\begin{frame}[shrink=25]{Zero initialization, and widening as you downsample}
\begin{minipage}[t]{0.494\textwidth}
\begin{icardaccent}
\cardh{kernel\_initializer="zeros"}
The last layer predicts \textbf{only zeros} after initialization --- so the model's default assumption before training is \textit{no noise}. A deliberately chosen, harmless starting point.
\end{icardaccent}
\end{minipage}%
\hfill
\begin{minipage}[t]{0.494\textwidth}
\begin{icardaccent}
\cardh{widths=[32, 64, 96, 128]}
Layers get \textbf{wider as the feature map shrinks}, and narrower again as it grows back. The same trade as every ConvNet in chapter 9.
\end{icardaccent}
\end{minipage}%
{\fontsize{9.4}{12.6}\selectfont\color{fg2}You would instantiate with \inlinecode{get\_model(image\_size=128, widths=[32, 64, 96, 128], block\_depth=2)}.\par}\vskip 4pt
\end{frame}

\framekicker{Section 17.2.3}
\begin{frame}[shrink=25]{Diffusion time and the diffusion schedule}
\begin{minipage}[t]{0.494\textwidth}
\begin{icardaccent}
\cardh{Diffusion time}
The index of the current step, here a \textbf{continuous value between 1 and 0}. 1 is maximal noise and minimal signal; 0 is almost all signal and no noise.
\end{icardaccent}
\end{minipage}%
\hfill
\begin{minipage}[t]{0.494\textwidth}
\begin{icardaccent}
\cardh{Diffusion schedule}
The relationship between the current time and how much noise and signal are present. We use a \textbf{cosine schedule}.
\end{icardaccent}
\end{minipage}%
{\fontsize{9.4}{12.6}\selectfont\color{fg2}The cosine choice is not arbitrary: it maintains the identity \hilite{noise\_rates$^{2}$ + signal\_rates$^{2}$ } 1==, so total energy stays constant as the mix shifts from one to the other.\par}\vskip 4pt
\end{frame}

\framekicker{Section 17.2.3 $\cdot$ listing 17.7}
\begin{frame}[shrink=25]{The schedule, in seven lines}
{\fontsize{9.4}{12.6}\selectfont\color{fg2}Map diffusion time to an angle, then take the cosine and sine of it. The Pythagorean identity does the rest.\par}\vskip 4pt
\icode{python}{listing 17.7}{listings/ch17-14.py}
{\fontsize{9.4}{12.6}\selectfont\color{fg2}The two bounds keep the process away from its extremes: never quite pure signal (0.95), never quite pure noise (0.02). Both endpoints are numerically awkward, and neither is needed.\par}\vskip 4pt
\end{frame}

\framekicker{Section 17.2.4}
\begin{frame}[shrink=25]{What the DiffusionModel needs to hold}
{\fontsize{9.4}{12.6}\selectfont\color{fg2}The denoising autoencoder is one attribute. Two more things are needed alongside it.\par}\vskip 4pt
\begin{ibullets}
\item \textbf{A loss function} --- mean absolute error, that is \inlinecode{mean(abs(real\_noise\_mask - predicted\_noise\_mask))}.
\item \textbf{An image normalization layer} --- the noise we add has unit variance and zero mean, so the images must be normalized the same way for their value ranges to match.
\end{ibullets}
\icode{python}{}{listings/ch17-15.py}
{\fontsize{9.4}{12.6}\selectfont\color{fg2}Mean absolute error rather than mean squared error: the noise mask is normally distributed, and \textbf{squared error would let a few extreme pixels dominate} the gradient.\par}\vskip 4pt
\end{frame}

\framekicker{Section 17.2.4}
\begin{frame}[shrink=25]{The denoise method}
{\fontsize{9.4}{12.6}\selectfont\color{fg2}Call the model for a predicted noise mask, then use it to reconstruct the clean image --- the one place the model's output is turned back into pixels.\par}\vskip 4pt
\icode{python}{}{listings/ch17-16.py}
{\fontsize{9.4}{12.6}\selectfont\color{fg2}Read the arithmetic as the inverse of the mixing formula: \inlinecode{noisy = signal * image + noise * mask}, so \inlinecode{image = (noisy $-$ noise * mask) / signal}. \textbf{Two lines that undo each other exactly.}\par}\vskip 4pt
\end{frame}

\framekicker{Section 17.2.4}
\begin{frame}[shrink=25]{The training step, in five operations}
\begin{center}
\includegraphics[width=0.92\linewidth,height=0.60\textheight,keepaspectratio]{figs/ch17-train-step.tex.pdf}
\end{center}
\vskip 3pt{\color{ink3}\fontsize{7.4}{9.4}\selectfont Random diffusion times are essential: the model must be trained across the full spectrum, because it will be called at every point of it.\par}
{\fontsize{9.4}{12.6}\selectfont\color{fg2}\inlinecode{call()} returns three things --- the predicted images, the predicted noise masks, and the \textbf{actual} noise masks it applied. The last two are what \inlinecode{compute\_loss()} compares.\par}\vskip 4pt
\end{frame}

\framekicker{Section 17.2.4}
\begin{frame}[shrink=25]{call() and compute\_loss(), written out}
{\fontsize{9.4}{12.6}\selectfont\color{fg2}The same pattern as the VAE: forward pass in \inlinecode{call()}, comparison in the loss.\par}\vskip 4pt
\icode{python}{}{listings/ch17-17.py}
{\fontsize{9.4}{12.6}\selectfont\color{fg2}Note how small the loss is. \textbf{All the difficulty is in the forward pass}, and the objective itself is one comparison.\par}\vskip 4pt
\end{frame}

\framekicker{Section 17.2.5}
\begin{frame}[shrink=25]{Generation: start from noise and walk the schedule down}
\begin{center}
\includegraphics[width=0.92\linewidth,height=0.60\textheight,keepaspectratio]{figs/ch17-generate.tex.pdf}
\end{center}
\vskip 3pt{\color{ink3}\fontsize{7.4}{9.4}\selectfont Each iteration denoises fully, then re-noises to the next schedule point --- a smaller step than the one just undone.\par}
{\fontsize{9.4}{12.6}\selectfont\color{fg2}That re-noising is the part that surprises people. The model predicts the \textbf{whole} clean image at every step; we then deliberately add back the amount of noise appropriate to the next time index, and ask again.\par}\vskip 4pt
\end{frame}

\framekicker{Section 17.2.5}
\begin{frame}[shrink=25]{The generation loop: one step}
{\fontsize{9.4}{12.6}\selectfont\color{fg2}Start from pure noise, then walk the schedule down one step at a time.\par}\vskip 4pt
\icode{python}{generate - the loop}{listings/ch17-18.py}
{\fontsize{9.4}{12.6}\selectfont\color{fg2}\inlinecode{diffusion\_steps} is a \textbf{generation-time} parameter --- the same weights sample in 5 steps or 50.\par}\vskip 4pt
\end{frame}

\framekicker{Section 17.2.5}
\begin{frame}[shrink=25]{The generation loop: re-noise and finish}
{\fontsize{9.4}{12.6}\selectfont\color{fg2}Having predicted the clean image, we deliberately add back the noise appropriate to the \textbf{next} time index.\par}\vskip 4pt
\icode{python}{generate - re-noise and denormalize}{listings/ch17-19.py}
{\fontsize{9.4}{12.6}\selectfont\color{fg2}The final two lines \textbf{undo the normalization} --- multiply by the standard deviation, add the mean, clip to [0, 255]. The model works in normalized space throughout; only the very last step returns to pixels.\par}\vskip 4pt
\end{frame}

\framekicker{Section 17.2.6}
\begin{frame}[shrink=25]{There is no metric, so look at the pictures}
{\fontsize{9.4}{12.6}\selectfont\color{fg2}We have no proper metric for image quality here. The practical answer is a callback that generates a grid at the end of every epoch so you can judge for yourself.\par}\vskip 4pt
\icode{python}{}{listings/ch17-20.py}
{\fontsize{9.4}{12.6}\selectfont\color{fg2}This is chapter 7's callback API doing something it was not obviously designed for, and doing it well.\par}\vskip 4pt
\end{frame}

\framekicker{Section 17.2.7}
\begin{frame}[shrink=25]{Two optimizer options that matter here}
{\fontsize{9.4}{12.6}\selectfont\color{fg2}We use \textbf{AdamW}, with two settings turned on to stabilise training and improve image quality.\par}\vskip 4pt
\begin{minipage}[t]{0.494\textwidth}
\begin{icardaccent}
\cardh{Learning rate decay}
An \inlinecode{InverseTimeDecay} schedule gradually reduces the rate through training.
\end{icardaccent}
\end{minipage}%
\hfill
\begin{minipage}[t]{0.494\textwidth}
\begin{icardaccent}
\cardh{Exponential moving average}
Also called \textbf{Polyak averaging}. Keep a running average of the weights, and every 100 batches overwrite the weights with it. Helps when the loss landscape is noisy.
\end{icardaccent}
\end{minipage}%
\icode{python}{}{listings/ch17-21.py}
{\fontsize{9.4}{12.6}\selectfont\color{fg2}Neither option changes what the model can represent. Both change \textbf{which minimum it settles into}, which for a generative model is the difference between plausible flowers and coloured smears.\par}\vskip 4pt
\end{frame}

\framekicker{Section 17.2.7}
\begin{frame}[shrink=25]{Fit it, and do not forget to adapt the normalizer}
{\fontsize{9.4}{12.6}\selectfont\color{fg2}Two lines to instantiate, one to train --- and one step in between that is easy to skip.\par}\vskip 4pt
\icode{python}{}{listings/ch17-22.py}
\begin{iband}[rose]
\inlinecode{model.normalizer.adapt(dataset)} computes the mean and variance the normalization needs. \textbf{Forget it and the noise and the images live on different scales}, and nothing works --- with no error message.
\end{iband}
{\fontsize{9.4}{12.6}\selectfont\color{fg2}100 epochs is about \textbf{90 minutes on a free Colab T4}.\par}\vskip 4pt
\end{frame}

\framekicker{Section 17.2.7 $\cdot$ figure 17.12}
\begin{frame}[shrink=25]{Flowers that do not exist}
\leadpar{After 100 epochs on 8,189 photographs, the model produces \textbf{convincing flowers}, and keeps improving if you keep training.}
{\fontsize{9.4}{12.6}\selectfont\color{fg2}It is worth measuring what that took against what it produced: a U-Net of a few million parameters, a dataset small enough to download over coffee, and ninety minutes of a free GPU.\par}\vskip 4pt
{\fontsize{9.4}{12.6}\selectfont\color{fg2}The next step to unlocking the real potential is \textbf{text conditioning} --- which turns this into a model that produces images matching a given caption.\par}\vskip 4pt
\end{frame}

\coursesection{04}{Text-to-image models}{One extra input to the denoiser, and a walk through the space between two prompts.}

\framekicker{Section 17.3}
\begin{frame}[shrink=25]{Give the denoiser a second input}
{\fontsize{9.4}{12.6}\selectfont\color{fg2}Take a pretrained \textbf{text encoder} --- a Transformer encoder like RoBERTa from chapter 15 --- that maps text to vectors in a continuous space. Then train a diffusion model on \textbf{(prompt, image)} pairs.\par}\vskip 4pt
\begin{center}
\includegraphics[width=0.92\linewidth,height=0.60\textheight,keepaspectratio]{figs/ch17-t2i-train.tex.pdf}
\end{center}
\vskip 3pt{\color{ink3}\fontsize{7.4}{9.4}\selectfont The denoising model now takes noisy\_images and text\_embeddings.\par}
{\fontsize{9.4}{12.6}\selectfont\color{fg2}This gives a considerable \textbf{leg up} on the flower denoiser: instead of removing noise with no additional information, the model gets a textual representation of the final image to guide it.\par}\vskip 4pt
\end{frame}

\framekicker{Section 17.3 $\cdot$ listing 17.8}
\begin{frame}[shrink=25]{Stable Diffusion in a few lines}
{\fontsize{9.4}{12.6}\selectfont\color{fg2}We will not train one from scratch --- you have all the ingredients, but it is expensive and slow. We use a pretrained model from KerasHub instead.\par}\vskip 4pt
\icode{python}{listing 17.8}{listings/ch17-23.py}
{\fontsize{9.4}{12.6}\selectfont\color{fg2}Like \inlinecode{CausalLM} last chapter, \inlinecode{TextToImage} is a \textbf{high-level task class} --- it wraps tokenization and the whole diffusion process into one \inlinecode{generate()} call. \inlinecode{dtype="float16"} is a memory trick explained properly in chapter 18.\par}\vskip 4pt
\end{frame}

\framekicker{Section 17.3}
\begin{frame}[shrink=25]{What that one call is doing}
{\fontsize{9.4}{12.6}\selectfont\color{fg2}Three stages hide behind \inlinecode{generate()}, and we will need all three separately later in this section.\par}\vskip 4pt
\begin{center}
\includegraphics[width=0.92\linewidth,height=0.60\textheight,keepaspectratio]{figs/ch17-t2i-steps.tex.pdf}
\end{center}
\vskip 3pt{\color{ink3}\fontsize{7.4}{9.4}\selectfont Encode the prompt, denoise conditioned on it, rescale the output.\par}
{\fontsize{9.4}{12.6}\selectfont\color{fg2}The middle stage is \textbf{exactly the flower model} --- the same loop over a diffusion schedule, with one extra input to the denoiser. Everything the previous section built is still here.\par}\vskip 4pt
\end{frame}

\framekicker{Section 17.3 $\cdot$ negative prompts}
\begin{frame}[shrink=25]{Steering away from something}
{\fontsize{9.4}{12.6}\selectfont\color{fg2}Stable Diffusion adds a second \textbf{negative prompt} input, used to steer the process away from certain text. There is nothing magic about it.\par}\vskip 4pt
{\fontsize{9.4}{12.6}\selectfont\color{fg2}Train on \textbf{triplets} --- \inlinecode{(image, positive\_prompt, negative\_prompt)}, where the positive describes the image and the negative is a series of words that do not. Feed both embeddings to the denoiser and it learns to move toward one and away from the other.\par}\vskip 4pt
\icode{python}{}{listings/ch17-24.py}
{\fontsize{9.4}{12.6}\selectfont\color{fg2}The result is the same scene with the blue removed --- a control surface built entirely out of \textbf{how the training data was labelled}.\par}\vskip 4pt
\end{frame}

\framekicker{Section 17.3 $\cdot$ sidebar}
\begin{frame}[shrink=25]{About those duplicated tusks}
{\fontsize{9.4}{12.6}\selectfont\color{fg2}Look closely at the output and you will find visual artifacts --- in the book's example, an elephant with duplicated tusks. Two separate causes, and only one is fixable.\par}\vskip 4pt
\begin{minipage}[t]{0.494\textwidth}
\begin{icardwarn}
\cardh{Unavoidable}
Drawing a human in a space suit sitting on a paper elephant would require understanding \textbf{anatomy and physics} the model lacks. It interpolates from training data with no real understanding of the objects.
\end{icardwarn}
\end{minipage}%
\hfill
\begin{minipage}[t]{0.494\textwidth}
\begin{icardgood}
\cardh{Easily fixable}
We are using the \textbf{smallest} Stable Diffusion 3 release --- about 3 billion parameters. A 9-billion-parameter version produces substantially fewer artifacts. It is omitted only to keep the example accessible.
\end{icardgood}
\end{minipage}%
\end{frame}

\framekicker{Section 17.3 $\cdot$ figure 17.15}
\begin{frame}[shrink=25]{Watching the denoising happen}
{\fontsize{9.4}{12.6}\selectfont\color{fg2}\inlinecode{num\_steps} exposes the loop count from section 17.2 as a generation parameter. Render the same prompt at several settings and the process becomes visible.\par}\vskip 4pt
\icode{python}{}{listings/ch17-25.py}
{\fontsize{9.4}{12.6}\selectfont\color{fg2}At 5 steps the image is a blurred impression; by 25 it is sharp. \textbf{The same weights throughout} --- only the number of times they were applied changed.\par}\vskip 4pt
\end{frame}

\framekicker{Section 17.3.1 $\cdot$ listing 17.9}
\begin{frame}[shrink=25]{Taking generate() apart}
{\fontsize{9.4}{12.6}\selectfont\color{fg2}To walk the latent space we need the three stages separately: encode the prompt, denoise conditioned on it, rescale the output.\par}\vskip 4pt
\icode{python}{listing 17.9}{listings/ch17-26.py}
{\fontsize{9.4}{12.6}\selectfont\color{fg2}Note the latents are \textbf{1/8 the image size} per dimension: Stable Diffusion denoises in a compressed space and decodes to pixels only at the end.\par}\vskip 4pt
\end{frame}

\framekicker{Section 17.3.1}
\begin{frame}[shrink=25]{The embedding is four tensors, not one}
\ioutput{listings/ch17-27.txt}
{\fontsize{9.4}{12.6}\selectfont\color{fg2}Rather than passing only the final embedded vector, the Stable Diffusion authors pass \textbf{both} the final output vector and the last representation of the entire token sequence --- more information for the denoiser to work with. Twice over, for positive and negative:\par}\vskip 4pt
\begin{ibullets}
\item The positive prompt's \textbf{encoder sequence} --- (1, 154, 4096)
\item The negative prompt's \textbf{encoder sequence} --- (1, 154, 4096)
\item The positive prompt's \textbf{encoder vector} --- (1, 2048)
\item The negative prompt's \textbf{encoder vector} --- (1, 2048)
\end{ibullets}
\end{frame}

\framekicker{Section 17.3.1 $\cdot$ listing 17.10}
\begin{frame}[shrink=25]{Interpolating on a sphere, not a line}
{\fontsize{9.4}{12.6}\selectfont\color{fg2}To walk between two prompts we interpolate their embeddings --- with \textbf{slerp}, spherical linear interpolation, a function used in computer graphics for decades.\par}\vskip 4pt
\icode{python}{listing 17.10}{listings/ch17-28.py}
{\fontsize{9.4}{12.6}\selectfont\color{fg2}The maths is not the point. \textbf{The motivation is.}\par}\vskip 4pt
\end{frame}

\framekicker{Section 17.3.1 $\cdot$ figure 17.16}
\begin{frame}[shrink=25]{Why linear interpolation goes wrong}
{\fontsize{9.4}{12.6}\selectfont\color{fg2}Imagine the text manifold as a sphere and two prompts as points on it. Interpolating \textbf{linearly} between them lands you \textbf{inside} the sphere --- off its surface.\par}\vskip 4pt
\begin{center}
\includegraphics[width=0.92\linewidth,height=0.60\textheight,keepaspectratio]{figs/ch17-slerp.tex.pdf}
\end{center}
\vskip 3pt{\color{ink3}\fontsize{7.4}{9.4}\selectfont Figure 17.16 --- spherical interpolation keeps us close to the surface of the manifold.\par}
{\fontsize{9.4}{12.6}\selectfont\color{fg2}The manifold is of course not actually spherical. But it is a smooth surface of numbers all of roughly the same magnitude --- it is \textbf{sphere-like}, and interpolating as if on a sphere is a better approximation than interpolating as if on a line.\par}\vskip 4pt
\end{frame}

\framekicker{Section 17.3.1}
\begin{frame}[shrink=25]{Interpolating the embeddings}
{\fontsize{9.4}{12.6}\selectfont\color{fg2}Only the two \textbf{positive} tensors are interpolated; the negative ones stay fixed, since we are not using them.\par}\vskip 4pt
\icode{python}{}{listings/ch17-29.py}
{\fontsize{9.4}{12.6}\selectfont\color{fg2}Elements 1 and 3 --- the negative prompt's sequence and vector --- are passed through unchanged from the first embedding.\par}\vskip 4pt
\end{frame}

\framekicker{Section 17.3.1 $\cdot$ figure 17.17}
\begin{frame}[shrink=25]{A dog becomes something else entirely}
{\fontsize{9.4}{12.6}\selectfont\color{fg2}Two prompts sharing a setting and differing in their subject, with nine images generated along the path between them.\par}\vskip 4pt
\icode{python}{}{listings/ch17-30.py}
{\fontsize{9.4}{12.6}\selectfont\color{fg2}The walk runs from \textbf{0.5 to 0.6} out of a full [0, 1] range --- zoomed into the middle of the interpolation, right where the morph becomes visually obvious. Nine images across a \textbf{one-tenth} slice of the path.\par}\vskip 4pt
\end{frame}

\framekicker{Section 17.3.1 $\cdot$ the closing argument}
\begin{frame}[shrink=25]{Interpolation machines}
\begin{iquote}
\quotetext{This might feel like magic the first time you try it, but there's nothing magic about it --- \textbf{interpolation is fundamental to the way deep neural networks learn.}}\quotecite{Section 17.3.1}
\end{iquote}
{\fontsize{9.4}{12.6}\selectfont\color{fg2}This is the last substantive model in the book, and a deliberately chosen visual metaphor to end on. Deep neural networks are \textbf{interpolation machines}: they map complex, real-world probability distributions onto low-dimensional manifolds.\par}\vskip 4pt
\begin{iband}[signal]
We can exploit that fact even for input as complex as human language and output as complex as natural images --- which is \hilite{exactly the claim chapter 15 made about attention}, arrived at from the other end.
\end{iband}
\end{frame}
\note{Tie the two together explicitly. Chapter 15 argued the point algebraically about embedding spaces; this chapter shows it, frame by frame. Same claim, two proofs.}

\framekicker{Common failure modes}
\begin{frame}[shrink=25]{Four ways generative image work goes wrong}
\begin{minipage}[t]{0.494\textwidth}
\begin{icardbad}
\cardh{Forgetting normalizer.adapt()}
The noise has unit variance; unnormalized images do not. Training runs, loss moves, and the output is \textbf{noise}. No error is raised.
\end{icardbad}
\end{minipage}%
\hfill
\begin{minipage}[t]{0.494\textwidth}
\begin{icardbad}
\cardh{Resizing without crop\_to\_aspect\_ratio}
Distortion in the training set becomes distortion in every generated image, and it is very hard to diagnose after the fact.
\end{icardbad}
\end{minipage}%


\vskip 4pt

\begin{minipage}[t]{0.494\textwidth}
\begin{icardwarn}
\cardh{Linear interpolation of embeddings}
Lands off the manifold. The midpoints decode to mush rather than to intermediate concepts. Use \textbf{slerp}.
\end{icardwarn}
\end{minipage}%
\hfill
\begin{minipage}[t]{0.494\textwidth}
\begin{icardwarn}
\cardh{Treating enhancement as evidence}
Super-resolution invents detail from its priors. The moon-crater example is the standing warning: \textbf{never for forensics}.
\end{icardwarn}
\end{minipage}%
\end{frame}

\framekicker{Summary}
\begin{frame}[shrink=25]{What has to survive this chapter}
\begin{isteps}
\item \textbf{Image generation means learning a latent space} that captures statistical information about a dataset --- sample it, decode it, and you get images nobody has drawn.
\item \textbf{A classical autoencoder's latent space is not useful}; a VAE's is, because encoding to a distribution plus a KL term forces continuity.
\item \textbf{VAEs give structured, controllable latent spaces} --- good for editing, concept vectors, and reconstruction, if not for fidelity.
\item \textbf{Override compute\_loss(), not train\_step()}, when leaving supervised learning --- it keeps the code working on all three backends.
\end{isteps}
\end{frame}

\framekicker{Summary}
\begin{frame}[shrink=25]{What has to survive this chapter (2 of 2)}
\begin{isteps}
\item \textbf{Diffusion is a denoising autoencoder in a loop}, predicting the noise mask rather than the clean image.
\item \textbf{The schedule is the design}: diffusion time runs from 1 to 0, and a cosine relationship keeps signal$^{2}$ + noise$^{2}$ = 1 throughout.
\item \textbf{A U-Net with skip connections} is the denoiser --- the same architecture as chapter 11's segmentation model, given a second scalar input.
\item \textbf{Text-to-image is one extra input} to the denoiser. Everything else is the flower model.
\item \textbf{Interpolation is what these models do.} Slerp between two prompts and you can watch it happen, frame by frame.
\end{isteps}
\vskip 4pt
\reslink{NOTEBOOK}{05\_stable\_diffusion\_latent\_walk.ipynb}{../../course-slides/notebooks/ch17/05\_stable\_diffusion\_latent\_walk.ipynb}
\reslink{PAPER}{Kingma \& Welling, Auto-Encoding Variational Bayes}{https://arxiv.org/abs/1312.6114}
\reslink{NEXT}{Chapter 18 --- Best practices for the real world}{../ch18/index.html}
\vskip 2pt
\end{frame}

\end{document}
